\section{System Architecture and Operating Conditions}
\label{sec:system}
The target and follower are SunFounder PiCar-S vehicles. The target executes an open-loop motion sequence; its commands and pose are unavailable to the follower. A speaker supplies a continuous $4\,\mathrm{kHz}$ reference tone. On the follower, a MAX4466 microphone inside a directional horn is rotated by a servo. The ESP32 extracts tone amplitudes at two horn angles and transmits timestamped pairs over USB. A Raspberry Pi estimates the relative state, computes closed-loop feedback, and commands forward speed and front steering.

\textbf{Platform correction.} The differential-drive description in the retained literature review does not describe the implemented vehicle. The PiCar-S has the bicycle constraint $\omega_f=v_f\tan\delta_s/L$ and cannot rotate its chassis at zero forward speed. Acquisition therefore rotates only the horn. References to differential-drive control in the review provide context, not the actuation model used below. The review's uniqueness statement is a motivation for evaluating this particular assembly, not a demonstrated novelty claim.

\begin{figure}[!t]
\centering
\resizebox{\linewidth}{!}{%
\begin{tikzpicture}[
  font=\footnotesize,
  blk/.style={draw, rounded corners=3pt, align=center, minimum height=11mm, minimum width=34mm, inner sep=3pt},
  hw/.style={blk, fill=blue!8},
  sw/.style={blk, fill=orange!14},
  link/.style={-{Latex[length=2mm]}, thick},
  both/.style={{Latex[length=2mm]}-{Latex[length=2mm]}, thick},
  snd/.style={-{Latex[length=2mm]}, thick, dashed},
  lbl/.style={font=\scriptsize, align=center},
  box/.style={draw, dashed, rounded corners=4pt, gray}
]
% ---- rover containers
\draw[box] (-0.4,-1.1) rectangle (3.6,1.5);
\node[lbl, anchor=north west, gray!70!black] at (-0.3,1.45) {\textbf{target rover}};
\draw[box] (4.9,-4.9) rectangle (13.6,1.5);
\node[lbl, anchor=north west, gray!70!black] at (5.0,1.45) {\textbf{follower rover}};
% ---- components
\node[hw] (tgt) at (1.6,0) {Speaker tone $f_0$\\OL motion};
\node[hw] (mic) at (7.2,0) {Servo arm with\\horn + MAX4466 mic};
\node[sw] (esp) at (11.7,0) {ESP32\\ADC, amplitude\\features, servo PWM};
\node[sw] (pi)  at (11.7,-3.3) {Raspberry Pi\\estimation and\\control software};
\node[hw] (drv) at (7.2,-3.3) {Motor and steering\\drivers};

% ---- links
\draw[snd] (tgt.east) -- node[above, lbl] {sound} (mic.west);
\draw[link] ([yshift=3mm]mic.east) -- node[above, lbl] {analog} ([yshift=3mm]esp.west);
\draw[link] ([yshift=-3mm]esp.west) -- node[below, lbl] {PWM} ([yshift=-3mm]mic.east);
\draw[both] (esp.south) -- node[right, lbl] {USB serial} (pi.north);
\draw[link] (pi.west) -- node[above, lbl] {commands} (drv.east);
\end{tikzpicture}%
}
\caption{System architecture: hardware components and communication links. Blue: physical elements; orange: computation. The signals and the control loop are shown in Fig.~\ref{fig:ctrl}.}
\label{fig:arch}
\end{figure}




The implemented system is a continuous nonlinear plant with sampled estimation and control and a discrete supervisory state. A local linear time-invariant (LTI) model supports gain design; the nonlinear simulator retains independent global poses, steering and speed lag, acceleration limits, and sequential acoustic captures. The controller receives simulated sensor packets through the same estimation interface used for recorded data. It has no access to simulated target truth.

The operating domain is planar motion on a flat indoor floor, a single known source, fixed source level, and low vehicle speeds. Range refers to the receiver--speaker reference points; chassis clearance is a separate geometric quantity. Speaker orientation, reflections, motor noise, and within-pair motion can violate the sensing assumptions. These effects require measured bounds before the simulation gains are transferred to the vehicles.

\begin{table}[t]
\caption{Nominal simulation assumptions; these are not measured hardware specifications.}
\label{tab:assumptions}
\centering\footnotesize
\begin{tabularx}{\columnwidth}{@{}Xr@{}}\toprule
Quantity & Value\\\midrule
Desired separation $d^\ast$ & $1.20\,\mathrm{m}$\\
Nominal speed $v_0$ & $0.18\,\mathrm{m/s}$\\
Wheelbase $L$ & $0.145\,\mathrm{m}$\\
Forward speed limits & $0.06$--$0.35\,\mathrm{m/s}$\\
Maximum steering magnitude & $25^\circ$\\
Speed / steering time constants & $0.25$ / $0.15\,\mathrm{s}$\\
Braking deceleration $a_b$ & $0.35\,\mathrm{m/s^2}$\\
Reference-point clearance / margin & $0.40$ / $0.10\,\mathrm{m}$\\
Assumed calibrated range domain & $0.55$--$3.50\,\mathrm{m}$\\
Horn width $\sigma$ / half-dither $\delta$ & $40^\circ$ / $20^\circ$\\\bottomrule
\end{tabularx}
\end{table}

\begin{table}[t]
\caption{Timing configuration and unresolved physical timing.}
\label{tab:timing}
\centering\footnotesize
\begin{tabularx}{\columnwidth}{@{}Xr@{}}\toprule
Quantity & Value / status\\\midrule
Control interval $T_s$ & $0.05\,\mathrm{s}$\\
Nominal pair period & $0.40\,\mathrm{s}$\\
Requested within-pair gap & $0.18\,\mathrm{s}$\\
Realized nominal simulated gap & $0.20\,\mathrm{s}$\\
Configured transport delay & $0.04\,\mathrm{s}$\\
Maximum accepted measurement age & $0.60\,\mathrm{s}$\\
Loss timeout / replay history & $1.30$ / $2.50\,\mathrm{s}$\\
Legacy static-log median pair period & $12.84\,\mathrm{s}$\\
Physical dwell, latency, jitter & \textcolor{red}{To measure}\\\bottomrule
\end{tabularx}
\end{table}
The simulator schedules captures and delivery on the $T_s$ grid. Thus the requested $0.18\,\mathrm{s}$ gap becomes $0.20\,\mathrm{s}$, and transport delivery is also quantized. A $20\,\mathrm{Hz}$ controller does not create $20\,\mathrm{Hz}$ acoustic information. The legacy static log contains five pairs, none usable, and cannot establish the required measurement rate.

\pending{H1 --- Hardware configuration}{Record measured wheelbase, reference-point offsets, steering calibration, motor deadband, speed map, actuator lag, braking response, speaker settings, and the final wiring/power arrangement. Replace assumed values only with traceable measurements.}
